\section{Introduction}
\label{sec:intro}

Almost every deployed language model answers users through a system prompt. A tutoring app tells the model to be ``Professor Hart''; a customer-facing assistant includes a page of safety rules; an integration requires JSON output; a chat widget asks for ``under 50 words''. Each prompt is written for one purpose, but it stays in front of {\em every} request. If a safety paragraph quietly costs accuracy, or a brand persona doubles token spend on unrelated questions, the people who write these prompts should know.

This question also matters for evaluation. Benchmark numbers are usually reported with no system prompt or a minimal one, while users interact with the same model under a long, deployment-specific prompt. Whether the benchmark numbers carry over depends on how large these {\em residual effects} are.

Prior work studies parts of this question in isolation. Personas in system prompts give small and inconsistent changes in factual accuracy~\cite{zheng2023helpful}. Forcing structured output such as JSON can hurt reasoning~\cite{tam2024speak}. Safety-tuned models refuse some benign requests~\cite{rottger2023xstest,cui2024orbench}. Superficial prompt-format changes can move accuracy by large margins~\cite{sclar2023quantifying}. However, these studies use different tasks, models and baselines, so they cannot tell a practitioner which {\em kind} of system prompt is costly. They also rarely separate the content of a prompt from its length, or a model's correctness from whether its answer can be parsed.

{\bf Which system prompts leave residue?} We answer this with a controlled factorial experiment (\figref{fig:design}). We write three wording variants for each of four common categories (\safety, \persona, \format and \constraint) and add three controls: no system prompt, the minimal ``You are a helpful assistant.'' (our baseline), and a length-matched but behaviourally empty deployment blurb. We hold the user turn fixed and run every condition on five benchmarks and three open instruction-tuned models. We score each response twice: a lenient scorer that estimates whether the content is right, and a strict scorer that mimics a production parser. We also measure output length, hedging and refusals. All effects are paired item-level differences from the baseline.

We tested four hypotheses drawn from common practitioner intuitions: safety prompts cost 4--8 points through hedging and refusal (H1); personas lengthen outputs by 30--50\% without helping (H2); format prompts are the only helpful category (H3); and constraints have the most variable effects, helping some tasks and hurting others (H4). The data contradict most of these intuitions:
\begin{itemize}[leftmargin=*,itemsep=0pt,topsep=0pt]
    \item \safety prompts are nearly free: $+0.2$\pp overall, $+3.7$\pp on \tqa, with hedging up to $5\times$ the baseline rate but refusals at or below 1\%.
    \item \persona prompts roughly {\em double} output length ($1.99\times$, well above the predicted 1.3--1.5$\times$) and cost 3.1\pp of accuracy.
    \item \format prompts reduce content accuracy ($-2.5$\pp) but improve strict-parser accuracy by $+19.6$\pp; whether they ``help'' depends entirely on the metric.
    \item \constraint prompts are the most harmful category ($-9.0$\pp) and never help: their variance comes from ``hurts a little'' \versus ``hurts badly'' (down to $-61$\pp when a model is told not to show its work), not from helping some tasks.
\end{itemize}

\noindent Our contributions are:
\begin{itemize}[leftmargin=*,itemsep=0pt,topsep=0pt]
    \item We design a controlled, side-by-side comparison of four system-prompt categories with length-matched controls, fixed user turns and paired item-level statistics over 135,000 generations.
    \item We separate content accuracy from parseability with two scorers, and show that this choice flips the sign of the conclusion for format prompts ($-2.5$\pp \versus $+19.6$\pp).
    \item We trace how items flip from right to wrong under each category, identifying concrete failure modes: compressed reasoning under schemas, rule spill-over across task types, and chain-of-thought removal under brevity limits.
    \item We show that the results are robust to wording variants, a length-matched control, answered-only scoring and a full temperature-0.7 replicate, and that residual effects depend strongly on the model.
\end{itemize}

The rest of the paper is organized as follows. \Secref{sec:related} reviews related work, \secref{sec:method} describes the design, \secref{sec:results} reports results, \secref{sec:discussion} discusses mechanisms and limitations, and \secref{sec:conclusion} concludes.

\begin{figure}[t]
\centering
\resizebox{0.95\linewidth}{!}{%
\begin{tikzpicture}[
    font=\small,
    box/.style={draw, rounded corners=2pt, align=center, minimum height=1.0cm, inner sep=4pt},
    arr/.style={-{Stealth[length=2mm]}, thick}]
  \node[box, fill=humanpurple, text width=4.3cm] (sys) at (0,1.0) {{\bf System prompt} (15)\\[1pt]\scriptsize controls: \nosys, \baseline, \neutrallong\\ \safety, \persona,\\ \format, \constraint\\ $\times$ 3 wording variants each};
  \node[box, fill=aigreen, text width=4.3cm] (usr) at (0,-1.0) {{\bf Fixed user turn}\\[1pt]\scriptsize minimal instruction + item\\ 5 tasks $\times$ 300 items};
  \node[box, fill=aigold!40, text width=3.0cm] (mod) at (4.6,0) {{\bf Model} (3)\\[1pt]\scriptsize \llama\\ \qwentwo\\ \qwenthree\\ greedy and $T{=}0.7$};
  \node[box, fill=white, text width=3.4cm] (score) at (8.6,0) {{\bf Score each response}\\[1pt]\scriptsize lenient \& strict accuracy\\ output tokens\\ hedge / refusal / unparseable};
  \node[box, fill=aired!50, text width=3.4cm] (stat) at (12.8,0) {{\bf Paired effect}\\[1pt]\scriptsize $\Delta$ \versus \baseline per item\\ bootstrap CI\\ McNemar test, Holm correction};
  \draw[arr] (sys.east) -- (mod.west);
  \draw[arr] (usr.east) -- (mod.west);
  \draw[arr] (mod) -- (score);
  \draw[arr] (score) -- (stat);
\end{tikzpicture}}
\caption{Experimental design. Only the system prompt changes between conditions; the user turn, decoding settings and scoring are identical. Every effect is a paired, item-level difference from the minimal baseline ``You are a helpful assistant.'' (\baseline), computed within each model.}
\label{fig:design}
\end{figure}
